% GENERATED FILE. Edit canonical lessons, then run npm run generate:books.
% canonical-source-hash: fnv1a64:b4e1db6c0499c4af
% canonical-lessons: JA-C87-tsumu, JA-C87-tanomu, JA-C87-susumu, JA-C87-hosu, JA-C87-sasu

\chapter{To pile up, To ask for, To go forward, To hang out to dry, To point}
\label{ch:ja-to-pile-up-to-ask-for-to-go-forward-to-hang-out-to-dry-to-point}
\hlchaptermodality{87}

\begin{chapteropening}
\textbf{By the end of this chapter:} \emph{I can say to pile up, to ask for, to go forward, to hang out to dry and to point.}

Complete the last lesson of chapter 87: i can say to pile up, to ask for, to go forward, to hang out to dry and to point.
\end{chapteropening}

\section[\texorpdfstring{\ja{つむ} (tsumu)}{つむ (tsumu)}]{\ja{つむ} (tsumu) — to pile up}

\label{lesson:JA-C87-tsumu}



\begin{quote}
Before the new one: say the Japanese for to step on, then the Japanese for to bite.
\end{quote}

\subsection*{You'll want to know: \ja{つむ}}
\textbf{\ja{つむ}} — \emph{tsumu} — \textquotedblleft{}to pile up\textquotedblright{}.

\begin{grammarlens}[title={What you've built}]
One more word for this chapter.
\end{grammarlens}

\subsection*{Your turn}
\begin{itemize}
\raggedright
  \item \emph{Say it:} \emph{tsumu}
  \item \emph{Say it:} \emph{tsumu}, once more
  \item \emph{Say it:} say \emph{tsumu}
  \item \emph{Recall:} say the Japanese for to step on, then the Japanese for to bite, then say \emph{tsumu} again
\end{itemize}

\begin{tcolorbox}[breakable,colback=teal!4,colframe=teal!35!black,title={Before you move on}]
What does \ja{つむ} mean? (\textquotedblleft{}To pile up\textquotedblright{}.) Say it once more.
\end{tcolorbox}
\section[\texorpdfstring{\ja{たのむ} (tanomu)}{たのむ (tanomu)}]{\ja{たのむ} (tanomu) — to ask for}

\label{lesson:JA-C87-tanomu}



\begin{quote}
Before the new one: say the Japanese for to bite, then the Japanese for to pile up.
\end{quote}

\subsection*{You'll want to know: \ja{たのむ}}
\textbf{\ja{たのむ}} — \emph{tanomu} — \textquotedblleft{}to ask for\textquotedblright{}.

\begin{grammarlens}[title={What you've built}]
One more word for this chapter.
\end{grammarlens}

\subsection*{Your turn}
\begin{itemize}
\raggedright
  \item \emph{Say it:} \emph{tanomu}
  \item \emph{Say it:} \emph{tanomu}, once more
  \item \emph{Say it:} say \emph{tanomu}
  \item \emph{Recall:} say the Japanese for to bite, then the Japanese for to pile up, then say \emph{tanomu} again
\end{itemize}

\begin{tcolorbox}[breakable,colback=teal!4,colframe=teal!35!black,title={Before you move on}]
What does \ja{たのむ} mean? (\textquotedblleft{}To ask for\textquotedblright{}.) Say it once more.
\end{tcolorbox}
\section[\texorpdfstring{\ja{すすむ} (susumu)}{すすむ (susumu)}]{\ja{すすむ} (susumu) — to go forward}

\label{lesson:JA-C87-susumu}



\begin{quote}
Before the new one: say the Japanese for to pile up, then the Japanese for to ask for.
\end{quote}

\subsection*{You'll want to know: \ja{すすむ}}
\textbf{\ja{すすむ}} — \emph{susumu} — \textquotedblleft{}to go forward\textquotedblright{}.

\begin{grammarlens}[title={What you've built}]
One more word for this chapter.
\end{grammarlens}

\subsection*{Your turn}
\begin{itemize}
\raggedright
  \item \emph{Say it:} \emph{susumu}
  \item \emph{Say it:} \emph{susumu}, once more
  \item \emph{Say it:} say \emph{susumu}
  \item \emph{Recall:} say the Japanese for to pile up, then the Japanese for to ask for, then say \emph{susumu} again
\end{itemize}

\begin{tcolorbox}[breakable,colback=teal!4,colframe=teal!35!black,title={Before you move on}]
What does \ja{すすむ} mean? (\textquotedblleft{}To go forward\textquotedblright{}.) Say it once more.
\end{tcolorbox}
\section[\texorpdfstring{\ja{ほす} (hosu)}{ほす (hosu)}]{\ja{ほす} (hosu) — to hang out to dry}

\label{lesson:JA-C87-hosu}



\begin{quote}
Before the new one: say the Japanese for to ask for, then the Japanese for to go forward.
\end{quote}

\subsection*{You'll want to know: \ja{ほす}}
\textbf{\ja{ほす}} — \emph{hosu} — \textquotedblleft{}to hang out to dry\textquotedblright{}.

\begin{grammarlens}[title={What you've built}]
One more word for this chapter.
\end{grammarlens}

\subsection*{Your turn}
\begin{itemize}
\raggedright
  \item \emph{Say it:} \emph{hosu}
  \item \emph{Say it:} \emph{hosu}, once more
  \item \emph{Say it:} say \emph{hosu}
  \item \emph{Recall:} say the Japanese for to ask for, then the Japanese for to go forward, then say \emph{hosu} again
\end{itemize}

\begin{tcolorbox}[breakable,colback=teal!4,colframe=teal!35!black,title={Before you move on}]
What does \ja{ほす} mean? (\textquotedblleft{}To hang out to dry\textquotedblright{}.) Say it once more.
\end{tcolorbox}
\section[\texorpdfstring{\ja{さす} (sasu)}{さす (sasu)}]{\ja{さす} (sasu) — to point; to hold up (an umbrella)}

\label{lesson:JA-C87-sasu}



\begin{quote}
Before the new one: say the Japanese for to go forward, then the Japanese for to hang out to dry.
\end{quote}

\subsection*{You'll want to know: \ja{さす}}
\textbf{\ja{さす}} — \emph{sasu} — \textquotedblleft{}to point; to hold up (an umbrella)\textquotedblright{}.

\begin{grammarlens}[title={What you've built}]
One more word for this chapter.
\end{grammarlens}

\subsection*{Your turn}
\begin{itemize}
\raggedright
  \item \emph{Say it:} \emph{sasu}
  \item \emph{Say it:} \emph{sasu}, once more
  \item \emph{Say it:} say \emph{sasu}
  \item \emph{Recall:} say the Japanese for to go forward, then the Japanese for to hang out to dry, then say \emph{sasu} again
\end{itemize}

\begin{tcolorbox}[breakable,colback=teal!4,colframe=teal!35!black,title={Before you move on}]
What does \ja{さす} mean? (\textquotedblleft{}To point; to hold up (an umbrella)\textquotedblright{}.) Say it once more.
\end{tcolorbox}
